% !TITLE 长相思·山一程
% !AUTHOR 纳兰性德
% !DYNASTY 明清
% !GENRE 词
% !ORDER 5

\begin{poem}
山一程\textsuperscript{1}，水一程，身向榆关\textsuperscript{2}那畔\textsuperscript{3}行，夜深千帐灯\textsuperscript{4}。

风一更\textsuperscript{5}，雪一更，聒碎乡心\textsuperscript{6}梦不成，故园\textsuperscript{7}无此声。
\end{poem}

\begin{annotations}
\notetext{1}{一程：一段路程。重复“一程”写跋山涉水、路途漫漫。}
\notetext{2}{榆关：即山海关，在今河北秦皇岛。清代皇帝常出山海关巡视东北，纳兰性德作为侍卫随行。}
\notetext{3}{那畔：那边、那一侧。指山海关以外的关外地区。}
\notetext{4}{千帐灯：千万座营帐中亮着灯火。行军途中夜晚宿营的壮观景象——但灯火虽多，没有一盏是家的灯。}
\notetext{5}{一更：古代夜间计时单位，一夜分为五更，每更约两小时。风一更、雪一更，写风雪整夜不停。}
\notetext{6}{聒碎乡心：风雪声喧扰，把思乡的心撕碎了。聒（guō），声音嘈杂扰人。}
\notetext{7}{故园：家乡、故里。家乡的夜晚没有这种风雪之声——言外之意，故园是安静的、温暖的。}
\end{annotations}

\begin{appreciation}
这首小令是纳兰性德随康熙皇帝出山海关巡视时写的，写旅途中的思乡——中国词里写思乡最干净的一首之一。

"山一程，水一程"，一段路，又一段路。走了多远？不知道，反正山完了是水，水完了是山。这让人想起《行行重行行》的开头："行行重行行，与君生别离。"两千年前的征人也是这么走的——中国的诗人，走了三千年，走的还是同一条路。

"身向榆关那畔行"，榆关就是山海关，往关外去。"那畔"是口语——那边。纳兰是满洲贵族，祖辈的天地本来就在关外，可他的家早就在北京城了。人往家相反的方向走，一步远过一步。

"夜深千帐灯"——千万座营帐灯火通明，是皇家出行的排场，壮观得很。可这五个字底下藏着另一层意思：灯再多，没有一盏是家里的灯。把灯火写得越亮，孤寂越深。这是这首词最厉害的地方。

"风一更，雪一更"，更是一夜五更的计时——风一整夜，雪一整夜。"聒碎乡心梦不成"，"聒"是吵——风雪声把想家的心吵碎了，觉也睡不成。末句是答案："故园无此声。"家乡的夜里没有这种声音。原来他睡不着，不只是因为风雪吵，是因为这声音太陌生——家是安静的、暖和的、熟悉的地方。

思乡是《诗经》就有的题目。《采薇》里写："曰归曰归，岁亦莫止。"想回家，想到一年到头回不去。三千年里，这首诗的意思没有变过。纳兰特别的地方在于，他不是被抓壮丁的穷人，是御前侍卫，前呼后拥，风光无限——可他在千帐灯火里想的，还是家。想家这件事，跟地位没有关系。

你以后第一次住校，半夜醒来听见陌生的声音，想家是很正常的事——一千年前，纳兰性德也在山海关外想他的家。别硬撑着，想家就给家里打个电话，妈妈的手机一直开着。
\end{appreciation}
